\abstract{РЕФЕРАТ}

Объем работы равен \formbytotal{lastpage}{страниц}{е}{ам}{ам}. Работа содержит 13 иллюстраций, 5 таблиц, \arabic{bibcount} библиографических источников и \formbytotal{числоПлакатов}{лист}{}{а}{ов} графического материала. Количество приложений – 2. Графический материал представлен в приложении А. Фрагменты исходного кода представлены в приложении Б.

Перечень ключевых слов: информационная система, 1С:Предприятие, конфигурация, спортивная организация, спортивная школа, справочник, документ, регистр, расписание занятий, соревнования, спортивные разряды, спортивный инвентарь, планировщик, управляемый интерфейс, ролевая модель, автоматизация учета.

Объектом разработки является информационная система на платформе 1С:Предприятие 8.3 для комплексного учета деятельности спортивной организации.

Целью выпускной квалификационной работы является разработка конфигурации, обеспечивающей автоматизацию ключевых бизнес-процессов спортивной организации: учет спортсменов и тренеров, составление расписания занятий, организацию соревнований, учет результатов и присвоение разрядов, складской учет спортивного инвентаря, формирование отчетности и разграничение прав доступа пользователей.

В процессе разработки на основе анализа предметной области была спроектирована структура конфигурации. Реализован алгоритм заполнения планировщика расписания с трехуровневой иерархической группировкой по видам спорта, тренерам и группам. Реализована функция импорта результатов соревнований из файлов Microsoft Excel. Настроена ролевая модель с ролями «Администратор» и «Тренер». Проведено функциональное тестирование системы.

Система функционирует в файловом режиме, что позволяет развернуть ее на одном компьютере без использования выделенного сервера. Разработанная конфигурация готова к внедрению в деятельность спортивных организаций.

\selectlanguage{english}
\abstract{ABSTRACT}

The volume of the work is 98 pages. The work contains 13 illustrations, 5 tables, 55 bibliographic sources, and 9 sheets of graphic material. The number of applications is 2. The graphic material is presented in Appendix A. The fragments of the source code are presented in Appendix B.

List of keywords: information system, 1C:Enterprise, configuration, sports organization, sports school, directory, document, register, class schedule, competitions, sports ranks, sports equipment, planner, managed interface, role model, account automation.

The object of development is an information system on the 1C:Enterprise 8.3 platform for comprehensive accounting of a sports organization's activities.

The purpose of the final qualification work is to develop a configuration that provides automation of key business processes in a sports organization: accounting for athletes and coaches, creating a training schedule, organizing competitions, recording results and assigning ranks, inventory accounting for sports equipment, generating reports, and differentiating user access rights.

During the development process, a configuration structure was designed based on an analysis of the subject area. An algorithm for filling the schedule planner with a three-level hierarchical grouping by sport, coach, and group was implemented. A function for importing competition results from Microsoft Excel files was implemented. A role model with "Administrator" and "Coach" roles was configured. Functional testing of the system was conducted.

The system operates in file mode, which allows it to be deployed on a single computer without using a dedicated server. The developed configuration is ready for implementation in the activities of sports organizations.

\selectlanguage{russian}